\documentclass[10pt,a4paper]{amsart}
\usepackage[margin=19mm]{geometry}
\usepackage{amsmath,amssymb,mathtools}
\usepackage[scr=rsfs,cal=euler]{mathalfa}
\usepackage{xfrac}
\usepackage[unicode,hidelinks,bookmarks=false]{hyperref}
\newcommand{\ie}{i.e.\ }
\newcommand{\cf}{cf.\ }
\newcommand{\resp}{respectively\ }
\newcommand{\Subsection}{\subsection}
\providecommand{\nobreakdash}{\nobreak\hbox{-}}
\setlength{\emergencystretch}{2em}
\multlinegap0pt
\usepackage{amssymb}
\usepackage{bm}
\usepackage{xcolor}
\newtheorem{lemma}{Lemma}
\newcommand{\dd}{\textup{d}}
\newcommand{\sE}{\mathscr{E}}
\newcommand{\sO}{\mathscr{O}}
\title{Transversely projectively flat foliations}
\author{St\'ephane Druel}
\date{Source: arXiv:2609.06642v1 (CC BY 4.0)}
\begin{document}
\maketitle

\noindent\emph{Source and licence.} Transversely projectively flat foliations. Version arXiv:2609.06642v1 (CC BY 4.0), distributed by arXiv under the Creative Commons Attribution 4.0 International licence, http://creativecommons.org/licenses/by/4.0/, as recorded in the arXiv licence metadata for this exact version and shown on the abstract page https://arxiv.org/abs/2609.06642v1. Full licence notice: \texttt{licenses/arxiv-2609.06642-CC-BY-4.0.txt}.\par\smallskip

\noindent\textit{Editorial scope.} Counted unit: the article's lemma lemma:proj\_flat\_frames (source0.tex lines 573--588). The article's complete original proof of it is delivered with it (source0.tex lines 590--610, one \texttt{proof} environment of 21 lines). No mathematics is written, repaired or re-derived: the statement and the proof are the article's own lines, in the article's own notation, and no retained line is substituted or re-flowed.  No further source-local context is needed: every cross-reference of every retained line resolves inside the two delivered spans.  Nothing was omitted from any retained span, so the package carries no omission record.  No retained line contains a citation, so the wrapper carries no bibliography and no bibliography entry.  No retained line contains a figure environment, an image-inclusion command or drawing code, so no image is delivered and the package declares no asset dependency.  Editorial formatting only, in addition: an AMS \texttt{amsart} preamble that restates the article's own declarations and macros used by the retained lines, namely the environments \texttt{lemma} on separate counters, together with the standard packages those lines need.  The retained environments print in this document as Lemma 1 in order of appearance, whereas the article prints the same items with the numbering of its own full document.  The complete native source of the article as distributed in the arXiv e-print for this exact version is bundled unchanged as \texttt{sources/209557/source0.tex}.
\begin{lemma}\label{lemma:proj_flat_frames}
The vector bundle $\sE$ is projectively flat if and only if there exists a covering $(U_i)_{i \in I}$ of $X$ by analytically open sets as well as frames $(e_1^i,\ldots,e_r^i)$ of $\sE$ over $U_i$ satisfying the following condition. There exist locally constant matrices $M_{ij} \in \textup{GL}(r,\sO_X(U_i))$ as well as nowhere vanishing holomorphic functions $f_{ij}$ on $U_i\cap U_j$ such that 
\[
\textbf{e}_i= f_{ij} M_{ij} \cdot \textbf{e}_j
\]
on $U_i \cap U_j$, where 
\[
\textbf{e}_i = 
\begin{pmatrix}
e_1^i \\
\vdots \\
e_r^i
\end{pmatrix}
.
\]
\end{lemma}

\begin{proof}
Suppose first that $\sE$ is projectively flat, and let $\square$ be a projective connection on $\sE$.
Let $(U_i)_{i \in I}$ be a covering of $X$ by analytically open sets, and let $\square_i$ be a flat connection on $\sE|_{U_i}$ lifting $\square|_{U_i}$. Shrinking the $U_i$, if necessary, we may assume that there exist frames
$(e_1^i,\ldots,e_r^i)$ of $\sE$ over $U_i$ such that $\square_i(\textbf{e}_i)=0$, where 
\[
\textbf{e}_i = 
\begin{pmatrix}
e_1^i \\
\vdots \\
e_r^i
\end{pmatrix}
.
\]
Let $\alpha_{ij}$ be closed holomorphic $1$-forms on $U_i\cap U_j$ such that $\square_i=\square_j+\alpha_{ij}\otimes \textup{id}$. Shrinking the $U_i$ further, we may assume that $\alpha_{ij}=\dd g_{ij}$ for some holomorphic function $g_{ij}$ on $U_i\cap U_j$. Let $N_{ij}\in \textup{GL}(r,\sO_X(U_i))$ be the transition matrix, so that $\textbf{e}_i=N_{ij} \cdot \textbf{e}_j$. We have
\begin{multline*}
0 = \square_i (\textbf{e}_i) = \square_i (N_{ij} \cdot \textbf{e}_j) = N_{ij} \cdot \square_i(\textbf{e}_j) + \dd N_{ij} \otimes \textbf{e}_j \\ = \alpha_{ij} \otimes \textbf{e}_i + \dd N_{ij}\cdot N_{ij}^{-1} \otimes \textbf{e}_i = \dd g_{ij} \otimes \textbf{e}_i + \dd N_{ij}\cdot N_{ij}^{-1} \otimes \textbf{e}_i.
\end{multline*}
This easily implies that $N_{ij}=\exp(-g_{ij}) M_{ij}$ for some locally constant matrix $M_{ij} \in \textup{GL}(r,\sO_X(U_i))$, so that $\textbf{e}_i= f_{ij} M_{ij} \cdot \textbf{e}_j$, where $f_{ij}=\exp(-g_{ij})$. 

Conversely, let $(U_i)_{i \in I}$ be a covering of $X$ by analytically open sets and let $(e_1^i,\ldots,e_r^i)$ be frames of $\sE$ over $U_i$ as in the statement of the lemma. Let $\square_i$ be the flat connection on $\sE|_{U_i}$ such that $\square_i (\textbf{e}_i)=0$. Then we have $\square_i=\square_j-\dd f_{ij}\otimes \textup{id}$ on $U_i \cap U_j$, and hence the $\square_i$ define a projective connection $\square$ on $\sE$.
\end{proof}

\end{document}
